\documentclass{article} % This command is used to set the type of
% document you are working on such as an article, book, or presenation
\usepackage{amssymb}
\usepackage{geometry} % This package allows the editing of the page layout
\usepackage{amsmath}  % This package allows the use of a large range
% of mathematical formula, commands, and symbols
\usepackage{graphicx}  % This package allows the importing of images

\newcommand{\question}[2][]{
  \begin{flushleft}
    \textbf{Question #1}: \textit{#2}

\end{flushleft}}
\newcommand{\sol}{\textbf{Solution}:} %Use if you want a boldface solution line
\newcommand{\maketitletwo}[2][]{
  \begin{center}
    \Large{\textbf{Assignment #1}

    Topology} % Name of course here
    \vspace{5pt}

    \normalsize{N. Ohayon Rozanes  % Your name here

    \today}        % Change to due date if preferred
    \vspace{15pt}

\end{center}}
\begin{document}
\maketitletwo[5]  % Optional argument is assignment number
%Keep a blank space between maketitletwo and \question[1]

\question[1]{}

We prove this by induction.
\begin{enumerate}
  \item $n = 0$, that is, the powerset of a set with no elements, or
    the empty set. Since the empty set is the subset of every set,
    including itself, $\mathcal{P}(\varnothing) = \{\varnothing\}$ is a set with
    one element, and therefore it has cardinality $1$. This verifies
    the hypothesis since $2^0 = 1$
  \item Assume this is true for $|X| = n$, that is, for $X$ being a set
    with $n$ elements $|\mathcal{P}(X)| = 2^n$
  \item For $|X| = n+1$, create $Y$ such that $Y \subset X$ and $|Y|
    = n$ so there is one element $x \in X, x \notin Y$. From $(2)$, we have that
    $|\mathcal{P}(Y)| = 2^n$. Now every subset of $X$ either contains
    $x$ or it does not. Those that do not are exactly the subsets of
    $Y$, and there are $2^n$ of them. Those that do are exactly the
    sets $S \cup \{x\}$ for $S \subset Y$, since this simply unions
    every subset, so there is clearly $2^n$ such subsets
    \[
      |\mathcal{P}(X)| = 2^n + 2^n = 2^{n+1}
    \]
    which is the hypothesis for $n+1$, completing the induction.
\end{enumerate}

\question[2]{}
%2
First, there is the trivial and gross topology:
\[
  \tau = \{\varnothing, X\} \quad \tau = \mathcal{P}(X)
\]
%5
There are three topologies which are constructed by adding a
singleton to the trivial topology.
\[
  \tau = \{\varnothing, \{a\}, X\}
\]
%8
Alternatively, there are three topologies where the smallest element is a pair
\[
  \tau = \{\varnothing, \{a,b\}, X\}
\]

%17
We can add any singleton to those, if the singleton is not in the
pair, then it's union with the pair would be the total set and its
intersection with the pair would be the empty set. If the singleton
is in the pair, it's intersection with the pair or total would be
itself and its union with the pair would be the pair. Topologies
constructed like this account for 9 additional topologies

\[
  \tau = \{\varnothing, \{a\}, \{a, b\}, X\}
\]
%20
For the topologies above with singletons not disjoint from the pair,
we can add the other singleton, since its union with the other
singleton would be the pair, and its intersection the empty set.
There are three such topologies
\[
  \tau = \{\varnothing, \{a\}, \{b\}, \{a,b\}, X\}
\]

%23
We could instead choose to include the other pair which intersects on
the singleton. There are three topologies like this

\[
  \tau = \{\varnothing, \{a\}, \{a,b\}, \{a, c\}, X\}
\]

The prior two families can be joined together to include two
singletons, their pair, and a pair with one of then and the disjoint
element, for example

\[
  \tau = \{\varnothing, \{a\}, \{b\}, \{a,b\}, \{a,c\}, X\}
\]

There are six such topologies, bringing the total to 29.

\question[3]{}

By construction, the topology $\tau$ includes the empty set. Further,
consider that $[-1, 1] \setminus [-1,1] = \varnothing$, the empty
set is finite of cardinality $0$, and thus is countable, therefore,
$[-1,1] \in \tau$.

Now we want to show that for any arbitrary collection of $U_i \in
\tau$, we have that

\[
  \bigcup_{i\in I} U_i \in \tau
\]
It suffices to show that
\[
  [-1,1] \setminus \bigcup U_i
\]
is countable. Consider that for any $k \in I$
\[
  U_k \subset \bigcup U_i \implies [-1,1]\setminus \bigcup U_i
  \subset [-1,1] \setminus U_k
\]

This is because if $x \in [-1,1] \setminus \bigcup U_i$, then $x \in
[-1,1]$ and $x \notin U_i$ for every $i$, in particular $x \notin U_k$, so
$x \in [-1,1] \setminus U_k$. By construction $[-1,1]\setminus U_k$
is countable, and a subset of a
countable set is countable, thus $[-1,1] \setminus \bigcup U_i$ is
countable and therefore in $\tau$

For finite intersection, notice that
\[
  [-1,1] \setminus \bigcap^n U_i \subset \bigcup [-1,1]\setminus U_i
\]

If $x \in [-1,1] \setminus \bigcap^n U_i$, then $x \in [-1,1]$
and $x \notin U_k$ for at least one $k$, so $x \in [-1,1] \setminus
U_k$, which is one of the sets being unioned on the right.

The finite union of countable sets is countable, and the subset of a
countable set is countable, therefore,
\[
  [-1,1] \setminus \bigcap U_i \in \tau
\]

Hence, we have verified that $\tau$ is a topology on $[-1,1]$

\question[4]{}
\begin{enumerate}
  \item $\forall x \in \mathbb{R}$, construct the set $B = [x, x+1)$.
    Then $x \in B \in \mathcal{B}_l$
  \item Let $B_1 = [a_1, b_1)$ and $B_2 = [a_2,
    b_2)$ be basis elements in $\mathcal{B}_l$ with $B_1 \cap B_2
    \neq \varnothing$
    then the intersection of those two elements
    \[
      B_1 \cap B_2 = [\max(a_1, a_2), \min(b_1, b_2)) \in \mathcal{B}_l
    \]
    As such, for any $x \in B_1 \cap B_2$, there
    exists a basis element $B_3 \subset B_1 \cap B_2$ such that $x \in B_3$
\end{enumerate}
It follows immediatly that $\mathcal{B}_l$ is a basis for $\mathbb{R}$

Similarily
\begin{enumerate}
  \item $\forall x \in \mathbb{R}$, construct the set $B = (x - 1,
    x]$, then $x \in B \in \mathcal{B}_u$
  \item Let $B_1 = (a_1 , b_1]$ and $B_2 = (a_2, b_2]$ be basis
    elements in $\mathcal{B}_u$ such that $B_1 \cap B_2 \neq
    \varnothing$. Then we have that
    \[
      B_1 \cap B_2 = (\max(a_1, a_2) \min(b_1, b_2)] \in \mathcal{B}_u
    \]
    As such, for any $x \in B_1 \cap B_2$ there exists a basis
    element, defined above, such that $B_3 \subset B_1 \cap B_2$ and $x \in B_3$
\end{enumerate}
Therefore, $\mathcal{B}_u$ is a basis for $\mathbb{R}$.

\question[5]{}
\begin{enumerate}
  \item For any $x \in \mathbb{R}, \exists$ an open interval $x \in (x-1,
    x+1) \in \mathcal{B}_k$
  \item Let $B_1, B_2 \in \mathcal{B}_k$ with $B_1 \cap B_2 \neq
    \varnothing$. We show that in fact $B_1 \cap B_2 \in
    \mathcal{B}_k$, so that we may simply take $B_3 = B_1 \cap B_2$
    for every $x \in B_1 \cap B_2$.

    Throughout, write $(a_1, b_1) \cap (a_2, b_2) = (c,d)$ where $c =
    \max(a_1,a_2)$ and $d = \min(b_1,b_2)$; this is again an open
    interval, and it is nonempty since $B_1 \cap B_2 \subset (a_1,b_1)
    \cap (a_2,b_2)$ is nonempty. There are three cases.
    \begin{itemize}
      \item If $B_1 = (a_1,b_1)$ and $B_2 = (a_2,b_2)$ are both open
        intervals, then $B_1 \cap B_2 = (c,d) \in \mathcal{B}_k$.
      \item If $B_1 = (a_1,b_1) \setminus K$ and $B_2 = (a_2,b_2)$,
        then
        \[
          B_1 \cap B_2 = \big((a_1,b_1) \setminus K\big) \cap (a_2,b_2)
          = \big((a_1,b_1) \cap (a_2,b_2)\big) \setminus K = (c,d)
          \setminus K \in \mathcal{B}_k
        \]
        where the middle equality is just the fact that removing $K$
        commutes with intersection: $x$ lies in the left side exactly
        when $x \in (a_1,b_1)$, $x \in (a_2,b_2)$, and $x \notin K$,
        which is exactly the condition for lying in the right side.
        The case $B_1 = (a_1,b_1)$, $B_2 = (a_2,b_2) \setminus K$ is
        identical with the roles reversed.
      \item If $B_1 = (a_1,b_1) \setminus K$ and $B_2 = (a_2,b_2)
        \setminus K$, then by the same reasoning
        \[
          B_1 \cap B_2 = \big((a_1,b_1) \cap (a_2,b_2)\big) \setminus
          K = (c,d) \setminus K \in \mathcal{B}_k
        \]
    \end{itemize}
    In every case $B_1 \cap B_2$ is itself an element of
    $\mathcal{B}_k$, so for any $x \in B_1 \cap B_2$ the choice $B_3 =
    B_1 \cap B_2$ satisfies $x \in B_3 \subset B_1 \cap B_2$.
\end{enumerate}
Therefore, $\mathcal{B}_k$ is a basis for $\mathbb{R}$.
\end{document}
